% 第二章 背景知识与相关工作
\section{背景知识与相关工作}

\subsection{密码学基础}

\subsubsection{Shamir 秘密共享}

Shamir 秘密共享（Shamir's Secret Sharing, SSS）\cite{shamir1979share} 是门限密码学的核心原语之一。其核心思想是将秘密 $S$ 拆分为 $n$ 个份额，使得任意 $t$ 个份额可以恢复完整秘密，而任意少于 $t$ 个份额则无法获得秘密的任何信息（信息论安全）。

\vspace{0.5em}
\noindent \textbf{数学原理}。在有限域 $\mathbb{Z}_p$ 上构造一个 $t-1$ 次多项式：

\[
f(x) = a_0 + a_1 x + a_2 x^2 + \cdots + a_{t-1} x^{t-1} \pmod{p}
\]

令 $a_0 = S$（主密钥），$a_1, \ldots, a_{t-1}$ 为随机选取的系数。计算 $n$ 个点 $(i, f(i))$ 作为份额分发给 $n$ 位监护人。恢复时使用\textbf{拉格朗日插值公式}：

\[
S = \sum_{i \in Q} y_i \cdot \prod_{j \in Q, j \neq i} \frac{-x_j}{x_i - x_j} \pmod{p}
\]

其中 $Q$ 为参与恢复的 $t$ 个份额的索引集合。

\vspace{0.5em}
\noindent \textbf{安全特性}。Shamir 方案具有完美安全性（perfect security）：任意 $t-1$ 个份额构成的方程组有 $p$ 个自由度，不能确定 $S$ 的任何信息。

\subsubsection{Pedersen 承诺}

Pedersen 承诺 \cite{pedersen1991non} 是一种信息论隐藏的计算绑定承诺方案。给定椭圆曲线 $\mathbb{E}$ 上的两个生成点 $G$ 和 $H$，满足 $H = G \cdot \text{SHA256}(\text{seed})$，即 $G$ 和 $H$ 的离散对数关系未知。

\vspace{0.5em}
\noindent \textbf{承诺生成}。对秘密值 $v$ 和盲因子 $r$，承诺为：

\[
C(v, r) = r \cdot G + v \cdot H
\]

\vspace{0.5em}
\noindent \textbf{安全属性}。

\begin{itemize}
    \item \textbf{完美隐藏}（Perfect Hiding）：对于任意 $C$ 和任意 $v$，存在 $r$ 使得 $C = rG + vH$，即 $C$ 不泄露 $v$ 的任何信息。
    \item \textbf{计算绑定}（Computational Binding）：找到 $(r', v') \neq (r, v)$ 使得 $C(r', v') = C(r, v)$ 等价于解椭圆曲线上的离散对数问题 $\log_G H$。
\end{itemize}

\vspace{0.5em}
\noindent \textbf{同态加法性质}。是本文方案的核心密码学基础：

\[
C(v_1, r_1) + C(v_2, r_2) = C(v_1 + v_2, r_1 + r_2)
\]

\subsubsection{椭圆曲线密码学}

本项目选用 \textbf{secp256k1} 椭圆曲线，该曲线由比特币社区广泛采用并经过充分的密码分析。secp256k1 的定义参数为：

\[
y^2 = x^3 + 7 \pmod{p}
\]
\[
p = 2^{256} - 2^{32} - 2^9 - 2^8 - 2^7 - 2^6 - 2^4 - 1
\]

曲线阶 $n \approx 2^{256}$，提供约 128 比特的安全强度。

\subsection{现有数字资产继承方案}

现有数字资产继承方案可大致分为以下几类：

\textbf{中心化托管方案}。以 Ledger Recover 为代表，用户私钥被加密分片后托管给第三方机构。此类方案的中心化信任模型与区块链去中心化的设计理念相悖，且需要 KYC 流程，用户的隐私仅限于"公司不会偷看"的信任假设。

\textbf{链上多签方案}。以 Safe\{Wallet\} 为代表，通过在智能合约中预设 m-of-n 签名策略控制资产转移。多签方案完全去中心化，但受限于区块链平台，仅支持链上资产，且所有交易公开可见，隐私性差。

\textbf{密码管理器继承方案}。通过紧急联系人功能，在用户长期不活动时将主密码发送给预设的接收人。此类方案本质上是"密码信封"模式，缺乏门槛机制——紧急联系人可获取完整密码，存在滥用和单点故障的风险。

值得一提的是，以上方案均未提供有效的\textbf{抗胁迫}设计。本项目的胁迫码机制填补了这一空白。


% ========== 参考文献条目（内联，供 bib 文件使用） ==========
%% 说明：以下仅为 bib 条目参考，实际引用在 ref.bib 中定义
%% [1] Shamir, A. How to share a secret. Communications of the ACM, 1979.
%% [2] Pedersen, T.P. Non-interactive and information-theoretic secure verifiable secret sharing. CRYPTO 1991.
